\documentclass[11pt,a4paper]{article}
\usepackage[utf8]{inputenc}
\usepackage[T1]{fontenc}
\usepackage[spanish]{babel}
\usepackage{hyperref}
\usepackage{geometry}
\geometry{margin=2.5cm}
\usepackage{parskip}

\title{341 las llaves y los candados son de la lucha libre}
\author{Alejandro Valenzuela Rivera\\\small Universidad Autónoma de Nuevo León}
\date{\today}
\begin{document}
\maketitle

\begin{quote}
\textbf{Sesgo declarado.} Este trabajo se sostiene desde el sesgo declarado de su autor, y no finge neutralidad. Lo que afirma se puede tumbar, y ese es el punto: el falsificador está escrito, no escondido.
\end{quote}

\section*{Resumen}
Las llaves y los candados vienen de la lucha libre: formas con nombre, que se pueden cerrar y abrir, y que otro puede aprender.

\section{Por qué se sostiene}
Una llave sin nombre no se enseña: se improvisa, y el que la recibe no puede repetirla.

\section{De dónde viene}
La autoridad: 'me dio la idea del nombre de llaves y candados'.

\section{El caso}
La lucha libre: formas con nombre que se cierran y se abren.

\section{Método}
\textquestiondown{}Otro puede repetir la llave, o solo quien la inventó?

\section{Límites, y qué lo tumba}
Una llave que solo funciona con su autor.

\section{Lo que queda abierto}
\textquestiondown{}La llave tiene nombre, y otro puede aprenderla?

\section*{Cómo se leyó esta tesis}
Esto se escribió en el repositorio man-made-empathy, del que este candado es una pieza: candados/06 del caso/341 las llaves y los candados son de la lucha libre.md. La máquina que lo compuso lo hizo desde este mismo archivo, y no agregó nada que no estuviera escrito. Si algo de esto es falso, se dice en el mismo sitio, y se corrige con fecha.

\end{document}
